
\section{Twentieth Sunday after Trinity: The Inheritance, the Fruit, and the Stone}

\begin{quote}

Matthew 21:33-44. Hear another parable: There was a householder, who planted a vineyard, and made a hedge around it, and dug a winepress in it, and built a tower, and let it out to husbandmen, and went into a far country. But when the time of the fruit drew near, he sent his servants to the husbandmen to receive its fruits. And the husbandmen took his servants; one they beat, another they killed, another they stoned. Again he sent other servants, more than the first, and they did likewise to them. But at last he sent his Son to them and said: They will surely have regard for my Son. But when the husbandmen saw the Son, they said among themselves: This is the Heir; come, let us kill him and seize upon his inheritance! And they took him and cast him out of the vineyard and killed him. When now the lord of the vineyard comes, what shall he do to those husbandmen? They said to him: He shall miserably destroy those wicked men, and let out the vineyard to other husbandmen, who shall give him the fruits in their seasons. Jesus said to them: Have you never read in the Scriptures: The stone which the builders rejected, the same has become the head of the corner; this is the Lord’s doing, and it is marvelous in our eyes? Therefore I say to you, the Kingdom of God shall be taken from you and given to a people that bears its fruits. And whoever falls upon this stone shall be broken to pieces; but the one upon whom it falls, him it shall crush.

\end{quote}

\bigskip

These are severe words of judgment that the Lord in this parable speaks to the people of Israel, and that because there were three things to which, in fleshly pride and vanity, they had not given heed.

The first is the inheritance.

All God’s promise to Israel turned upon the inheritance. At the Lord’s command Abraham had gone out from his home and his people into a strange land, to which otherwise he had no right, but concerning which God said to him: “To your seed will I give this land” (Gen. 12:7). Thus the land of Canaan became the inheritance of the people of Israel.

But this was only one part, and indeed only the outward and earthly part, of the inheritance which the Lord gave Abraham and his seed to possess through the covenant and the promises. Through it there was also promised the whole world, the glory and the giving of the Law and the worship of God and the Word, and above all things the filial election in the Son himself, the Messiah, the Savior of the world, the goal and end of all the promises, the Heir himself, to whom the Father had given the Gentiles for an inheritance (Ps. 2), and whose day Abraham saw and rejoiced.

Therefore the Lord had also said to Abraham: “In you shall all the families of the earth be blessed,” and Paul expresses it thus: that Abraham and his seed received the promise that they should inherit the world (Rom. 4:13). But he also adds: “Not through the Law, but through the righteousness of faith,” and still further he says (Rom. 4:16), “that the inheritance is by faith, in order that it may be by grace, and that the promise may stand fast.”

This was precisely what the people of Israel had forgotten. They thought that because in an outward and fleshly sense they were descended from Abraham, therefore they had the right of inheritance, although, as John and Jesus say, God can raise up children of Abraham of that sort from stones. They passed over faith, without which no one can become Abraham’s true spiritual seed, and grace, which abolishes all claims and without which the promises do not stand fast. They claimed the inheritance by virtue of descent and the outward works of the Law; but an inheritance cannot be taken, it must be given, and given only to the children, the true heirs. To be Abraham’s seed, Abraham’s faith is required; and to be God’s children, it is required that one be born of God through faith. A new heart is required, a heart of flesh in place of the heart of stone, and where this change has not taken place, there is no right of inheritance, because there is no life, and therefore no fruit for God.

That was the second thing the people of Israel forgot: to bear fruit for God.

In the fifth chapter of the prophet Isaiah this fruitlessness of the chosen people is portrayed in vivid strokes. What was it that the Lord had not done for his vineyard? And when he came to harvest good grapes, he found wild ones. What then shall he do with such a vineyard? “I will take away,” says the Lord, “its hedge, and it shall be devoured; I will break down its wall, and it shall be trodden down. And I will lay it waste; it shall not be pruned nor hoed, so that briers and thorns come up; and I will command the clouds that they rain no rain upon it; for the vineyard of the Lord of hosts is the house of Israel, and the men of Judah his beloved planting; and he looked for justice, and behold, bloodshed; for righteousness, and behold, a cry” (Isa. 5:5-7).

Year after year the Lord came to his fig tree to find fruit and found none (Luke 13:7). Year after year God’s servants cried to him and said: “Spare the tree yet this one year!” In vain they dug about it and manured it; the Lord’s beloved planting, the men of Judah, brought him no fruit; and when the Son himself came to bring the fruit home, the measure was full; on the fig tree there were only leaves and no fruit, and judgment had to come, as the Savior pronounces it in our parable, the judgment of scattering and destruction, the judgment of the curse: the Kingdom had to be taken from them and given to a people that would bear its fruit.

The people of Israel also gave no heed to the stone. It was a chief cornerstone and became to them a stone of stumbling.

Since they had no life, no fruit, and therefore no right of sonship, they wished to rob the inheritance and lay violent hands upon the Son, in whom alone are life and the status of children. They cast from them the blessed Word: “Behold, I have laid in Zion a stone, a tried stone, a precious, sure-founded cornerstone; whoever believes shall not hasten” (Isa. 28:16). And although God’s house had been entrusted to them for edification, they rejected this chief cornerstone and built the house of their self-righteousness upon sand; when then the torrents came, its fall was great. They were offended at this chief cornerstone, because he was lowly in appearance and yet judged their conscience and rejected their righteousness and offered grace only for repentance and a broken heart, that sacrifice and that fruit which alone please God. The prophet’s word was fulfilled upon them, so that he who is the true sanctuary “became a stone of stumbling and a rock of offense for both the houses of Israel, a snare and a trap for the inhabitants of Jerusalem” (Isa. 8:14).

And so it came to pass that when in blind rage they rushed forward against this living stone and hanged him upon the tree of the curse, their judgment was sealed, their foundation torn away from beneath them; and when he rose victorious from the dead and from the Father’s right hand sent forth his messengers with the Gospel to hand over the Kingdom to other peoples who would bear its fruit, then the stone fell upon the temple and crushed it, and upon the people and scattered it; the judgment was accomplished.

Brothers and sisters! If anyone is in Christ, then he is Abraham’s seed and, according to the promise, an heir; “as many as have been baptized into Christ have put on Christ.” You who have been marked and born by the holy water of baptism, you are now the heir; you are the one who shall bear the fruit of the Kingdom of God, you are set upon the chief cornerstone.

But if you do not bear fruit, then the inheritance shall be taken from you, and the stone which you have despised shall fall upon you and crush you. It shall go worse with you than with Israel; for if God did not spare the natural branches, neither shall he spare you. See therefore the goodness and severity of God: severity toward those who have fallen, but goodness toward you, if you continue in his goodness; otherwise you also shall be cut off (Rom. 11:21-22). Do not be proud, but fear!

“I say to you, unless you repent, you shall all likewise perish.”

Friend, where is the fruit of the Kingdom of God among us? Do you have fruit?
